\chapter{TRABALHOS RELACIONADOS}
\label{chap:trabalhos-relacionados}

\section{Integração de dados orbitais e geoquímicos em rejeitos de mineração}

\citeonline{zhang2023} desenvolveram um método para estimar a concentração de ouro ainda presente no depósito de rejeitos de uma mina de ouro em Witwatersrand, na África do Sul.

Os autores analisaram imagens dos satélites Landsat-8 e Sentinel-2 e utilizaram um modelo geoestatístico construído a partir de análises de amostras para representar os teores de ouro. Em seguida, realizaram a fusão de dados entre esse modelo e os dados espectrais do Sentinel-2, formando o conjunto utilizado na comparação de cinco algoritmos de aprendizado de máquina: kNN, SVM, Random Forest, AdaBoost e redes neurais artificiais. A seleção dos modelos e a avaliação das previsões foram realizadas por validação cruzada em 10 partes, utilizando o coeficiente de determinação ($R^2$), o erro absoluto mediano e o erro percentual absoluto médio. Os autores também relatam resultados de validação cruzada baseada em blocos, obtidos em 100 execuções.

Entre os algoritmos avaliados, AdaBoost apresentou os melhores resultados, com coeficiente de determinação ($R^2$) de 0,917 e erro percentual absoluto médio de 2,7\%. Os modelos permitiram gerar mapas de concentração de ouro com resolução espacial de 10 metros. Os autores destacam que a confiabilidade das previsões depende da semelhança entre os materiais presentes na área de aplicação e aqueles representados nos dados de treinamento.

O estudo se relaciona com o presente trabalho por utilizar dados espectrais de satélite para predizer uma variável geoquímica. A principal diferença está no modelo e nas tarefas investigadas: enquanto Zhang et al. avaliaram diferentes algoritmos para estimar o teor de ouro em rejeitos de mineração, este trabalho propõe avaliar o TabPFN na classificação litológica e na predição de composição mineral. Dessa forma, o artigo oferece um exemplo de integração entre dados espectrais e dados de referência que pode orientar a preparação dos dados e a análise dos resultados deste TCC.

\section{TabPFN em bibliotecas espectrais de solo}

\citeonline{barkov2026} investigaram modelos de regressão para estimar propriedades do solo a partir de espectros no visível e infravermelho próximo (vis-NIR) e no infravermelho médio (MIR). O objetivo foi avaliar o desempenho do TabPFN em diferentes escalas de dados e verificar se a redução de dimensionalidade poderia melhorar as previsões diante da grande quantidade de variáveis espectrais e da elevada correlação entre elas.

Os autores utilizaram duas coleções de dados, avaliadas separadamente. Da LimeSoDa, selecionaram 13 conjuntos com dados de espectroscopia, contendo entre 32 e 460 amostras de solo cada, e mantiveram apenas as variáveis espectrais quando havia informações de outros sensores. Cada conjunto foi utilizado para prever três propriedades determinadas por análises laboratoriais: carbono orgânico ou matéria orgânica, pH e teor de argila, totalizando 39 tarefas de regressão. Da Open Soil Spectral Library (OSSL), que reúne dados espectrais harmonizados de diferentes coleções, selecionaram as tarefas com pelo menos 10 mil observações da propriedade-alvo. Essa seleção resultou em 46 tarefas, sendo 30 com espectros MIR e 16 com vis-NIR, com conjuntos de 14.166 a 82.573 amostras. Os alvos incluíram propriedades como carbono, nitrogênio, pH, frações granulométricas e teor de carbonatos. Em ambas as coleções, os espectros constituíram as variáveis de entrada e as propriedades medidas nas amostras foram os valores de referência para a avaliação.

Foram comparados seis modelos: regressão linear, regressão por mínimos quadrados parciais (PLSR), Random Forest, Cubist, rede neural convolucional (CNN) e TabPFN, na versão 2.5. Também foi avaliado um ensemble de CNNs. Os autores combinaram modelos com espectros completos ou com variáveis obtidas por análise de componentes principais (PCA) e mínimos quadrados parciais (PLS), formando 14 configurações. Os espectros passaram por normalização e padronização, e a avaliação utilizou validação cruzada aninhada, com cinco partições externas aleatórias. Na etapa interna, empregaram validação cruzada em cinco partes para a LimeSoDa e uma divisão de validação com 20\% dos dados de treinamento para a OSSL. A raiz do erro quadrático médio (RMSE) foi utilizada para avaliar as previsões e ordenar as configurações em cada tarefa, permitindo comparar posições médias e medianas entre propriedades com escalas e unidades distintas.

O TabPFN apresentou o melhor desempenho global nas duas coleções, e sua combinação com variáveis latentes de PLS ocupou a primeira posição no ranking agregado de cada uma. Na OSSL, essa combinação obteve o menor RMSE em 45,7\% das tarefas e permaneceu a até 1\% do menor erro em 71,7\% delas. Esses percentuais descrevem a comparação entre configurações nas tarefas avaliadas, e não uma medida de acurácia. O Cubist combinado com PLS foi a alternativa clássica mais competitiva. Como limitações, os autores destacam o custo de inferência do TabPFN, que chegou a 24,1 minutos nas maiores tarefas e foi reduzido para até 9,6 minutos com PLS, e o uso de amostras de treinamento e teste provenientes da mesma população. Portanto, o estudo não demonstra que essa vantagem se mantém na aplicação a regiões, instrumentos ou condições de aquisição diferentes dos representados nos dados avaliados.

O estudo se relaciona com este TCC por avaliar o TabPFN na predição de propriedades a partir de dados espectrais organizados em formato tabular. Entretanto, seu foco é a regressão de propriedades do solo, sem avaliar classificação litológica. Embora os alvos incluam atributos relacionados à composição, como teor de argila e carbonatos, essas tarefas não equivalem à identificação de classes de rocha ou à predição das proporções de diferentes minerais. Este trabalho pretende investigar classificação litológica e predição da composição mineral a partir de dados espectrais de satélite, conforme a disponibilidade de rótulos e medidas de referência adequados a cada tarefa.

\citeonline{barkov2026} destacam que a agregação de dados de laboratórios com métodos e instrumentos distintos introduz variabilidade adicional, tratada na OSSL por procedimentos de harmonização. A diversidade dos conjuntos amplia a abrangência da avaliação dos modelos, mas o estudo não testa se a mistura de sensores, por si só, melhora a predição. Essa distinção é relevante para este TCC, que propõe integrar e harmonizar conjuntos de dados geológicos e geoquímicos de diferentes regiões do mundo. Pretende-se associar coordenadas, classes litológicas e medidas laboratoriais compatíveis aos valores de reflectância das imagens Sentinel-2, considerando a correspondência espacial e temporal entre as amostras e a superfície observada pelo satélite.

O projeto pretende investigar a hipótese de que a integração de dados harmonizados de diferentes regiões pode ampliar a representatividade do treinamento e melhorar a generalização do TabPFN. Para isso, propõe avaliar o modelo em regiões que não participaram do ajuste, comparando o uso de dados de uma região com o de múltiplas regiões. Também se pretende disponibilizar os conjuntos derivados no Kaggle, conforme as condições de redistribuição das fontes, preservando sua procedência e os critérios de compatibilidade. Essa proposta se distingue da avaliação com partições aleatórias realizada por Barkov et al. e aborda a transferência entre regiões, que os autores apontam como uma questão ainda aberta. O eventual ganho de desempenho constitui uma hipótese do TCC, a ser verificada experimentalmente.
